\section{Lecture 11 (October 14th)}
\begin{rmk}
A (long long long)$^3$ time ago, we learned about the existence of spliiting fields and isomorphism extensions. Today, we learn about how splitting fields are unique.
\end{rmk}
\vspace{2ex}
\begin{prop}
(14.1) Let $K\subset F$ be a field extension and $p(t)\in K[t]$ be an irreducible element. There is a one-to-one correspondence between the set 
\[\{\alpha \in F \;|\; p(\alpha )=0\}\]
and
\[\{\sigma :K[t]/(p(t))\rightarrow F\;|\;\sigma |_{K}=\mathrm{1}_{K}\}\]
\end{prop}
\vspace{2ex}
\begin{thm}
(Isomorphism extension for simple algebraic extensions)
\begin{center}
\begin{tikzpicture}[baseline=(current bounding box.center),
  >={Straight Barb[length=3pt,width=6pt]}, % barbed arrowhead (side lines)
  node font=\small]

  % corners
  \node (E) at (0,4.5) {$E$};
  \node (F) at (3,4.5) {$F$};
  \node (A) at (0,3) {$K(\alpha )$};
  \node (B) at (0,0) {$K$};
  \node (C) at (3,0) {$K$};
  \node (D) at (3,3) {$K(\beta )$};
  \node (G) at (-3,3) {$K[t]/(p(t))$};
  \node (H) at (6,3) {$K[t]/(p(t))$};

  % vertical arrows (pointing up)
  \draw[->] (B) -- (A);
  \draw[->] (C) -- (D);
  \draw[->] (A) -- (E);
  \draw[->] (D) -- (F);
  \draw[->] (A) --  node[midway,above] {$\cong$}(G);
  \draw[->] (D) --  node[midway,above] {$\cong$}(H); 
  \draw[->] (A) -- node[midway,above] {$\cong$}(D);
  \draw[->] (G) to[bend left=50] node[midway,above] {${\bf 1}$} (H);

  % horizontal arrows (pointing right) with labels above
  \draw[->] (B) -- node[midway, above] {$ {\bf 1}$} (C);

  % commutator symbol in the middle (clockwise)

\end{tikzpicture}
\end{center}
Let $\alpha \in E$ be algebraic over $K$ and $\beta \in F$ be a zero of $p(t)=\mathrm{irr}(\alpha ,K)$. Then, there exists an isomorphism $\sigma :K(\alpha )\rightarrow K(\beta )$ such that $\sigma |_{K}=\mathrm{1}_{K}$
\end{thm}
\vspace{2ex}
\begin{thm}
(Isomorphism extension for finite extensions) Let $E/K$ be a finite extension. Fix an algebraic closure $\bar{K}$ of $K$. Then there exists a $K$-embedding $\sigma :E\rightarrow \bar{K}$. 
\begin{center}
\begin{tikzpicture}[baseline=(current bounding box.center),
  >={Straight Barb[length=3pt,width=6pt]}, % barbed arrowhead (side lines)
  node font=\small]

  % corners
  \node (E) at (0,3) {$\big(K(\alpha_1)\big)(\alpha _2)$};
  \node (F) at (3,3) {$\big(K(\beta _1)\big)(\beta_2)$};
  \node (A) at (0,1.5) {$K(\alpha_1 )$};
  \node (B) at (0,0) {$K$};
  \node (C) at (3,0) {$K$};
  \node (D) at (3,1.5) {$K(\beta_1 )$};
  \node (G) at (0,4.5) {$E$};
  \node (H) at (3,4.5) {$\bar{K}$};

  % vertical arrows (pointing up)
  \draw[->] (B) -- (A);
  \draw[->] (C) -- (D);
  \draw[->] (A) -- (E);
  \draw[->] (D) -- (F);
  \draw[->] (A) -- node[midway,above] {$\sigma_1$}(D);
  \draw[->] (E) -- node[midway,above] {$\sigma_2$}(F);
  \draw[->] (E) -- (G);
  \draw[->] (F) -- (H);

  % horizontal arrows (pointing right) with labels above
  \draw[->] (B) -- node[midway, above] {$ {\bf 1}$} (C);

\end{tikzpicture}
\end{center}
\end{thm}
\vspace{2ex}
\begin{proof}
Since $E/K$ is finite, $E=K(\alpha_1,\alpha_2,\ldots ,\alpha_{n})$ for some algebraic elements $\alpha_1,\ldots ,\alpha_{n}$ over $K$. Using this fact and induction, we can prove our theorem.

\bigskip

For $\alpha \in E$ which is algebraic in $K$, we can set $p_1(t)=\mathrm{irr}(\alpha ,K)\in K[t]$. Let $\beta_1\in \bar{K}$, a zero of $p(t)$ in $\bar{K}$. Then, let $\alpha_2\in E$, which is again algebraic in $K(\alpha _1)$. We can set $p_2(t)=\mathrm{irr}(\alpha_2,K(\alpha_1))\in \big(K(\alpha_1)\big)[t]$. We can choose a zero $\beta_2\in  \bar{K}$ of $p_2(t)\in \big(K(\beta_1))\big)[t]$.  
\end{proof}
\vspace{2ex}
\begin{rmk}
Let $p(t)\in K[t]\subset \bar{K}[t]$ be an irreducible element and 
\[p(t)=(t-\alpha_1)^{m_1}(t-\alpha_2)^{m_2}\ldots (t-\alpha _{n})^{m_{n}}\]
is $m_{i}=m_{j}$ for $i\ne j$?
\end{rmk}
\vspace{2ex}

\begin{thm}
(Isomorphism extension theorem) Let $K$ be a field, and let $f(t)\in K[t]$ be a nonconstant polynomial. Suppose that $F_1/K$ and $F_2/K$ are splitting fields for $f(t)$. Then there exists a $K$-isomorphism $\sigma :F_1\rightarrow F_2$.
\end{thm}
\vspace{2ex}
\begin{proof}
\begin{itemize}
\item[(i)] Choose a zero $\alpha_1$ of $f(t)$ in $F_1$ and let $p(t)=\mathrm{irr}(\alpha_1,K)$
\item[(ii)] Choose a zero $\beta_1$ of $p(t)$ in $F_2$
\end{itemize}
We then have an isomorphism $\sigma_1$ between $K(\alpha_1)$ and $K(\beta _1)$. We can continue and achieve:
\begin{center}
\begin{tikzpicture}[baseline=(current bounding box.center),
  >={Straight Barb[length=3pt,width=6pt]}, % barbed arrowhead (side lines)
  node font=\small]

  % corners
  \node (E) at (0,3) {$\big(K(\alpha_1)\big)(\alpha _2)$};
  \node (F) at (3,3) {$\big(K(\beta _1)\big)(\beta_2)$};
  \node (A) at (0,1.5) {$K(\alpha_1 )$};
  \node (B) at (0,0) {$K$};
  \node (C) at (3,0) {$K$};
  \node (D) at (3,1.5) {$K(\beta_1 )$};
  \node (G) at (0,4.5) {$F_1$};
  \node (H) at (3,4.5) {$F_2$};

  % vertical arrows (pointing up)
  \draw[->] (B) -- (A);
  \draw[->] (C) -- (D);
  \draw[->] (A) -- (E);
  \draw[->] (D) -- (F);
  \draw[->] (A) -- node[midway,above] {$\sigma_1$}(D);
  \draw[->] (E) -- node[midway,above] {$\sigma_2$}(F);
  \draw[->] (E) -- (G);
  \draw[->] (F) -- (H);

  % horizontal arrows (pointing right) with labels above
  \draw[->] (B) -- node[midway, above] {$ {\bf 1}$} (C);

\end{tikzpicture}
\end{center}
\end{proof}
\vspace{2ex}
\begin{rmk}
Let $f(t)=(t-\alpha_1)\ldots (t-\lambda_{n})$ in $\bar{K}[t]$. Then there exists a $K$-isomorphism $\sigma _{i}:F_1\rightarrow K(\lambda_1,\ldots ,\lambda_{n})\subset \bar{K}$.
\end{rmk}
\vspace{2ex}

